\documentclass{article}
\usepackage{amsmath,amssymb,hyperref}
\title{Horizontal assembly of native state bicategories}
\date{15 November 2026}
\begin{document}\maketitle
\paragraph{Status.} Independently reviewed conditional written construction. The fixed-interface dialgebra bicategories and their synchronization are already checked. This note supplies the next higher-dimensional assembly data; it does not identify a tricategory merely from the existence of its hom bicategories.

\paragraph{Independent input.} Let $S$ be a symmetric monoidal bicategory of pure states, $K$ a process bicategory, and $J$ the specified pure pseudofunctor with coherent pure actions. Require strong left and right pure-action pseudofunctors, invertible action associativity/unit and left/right compatibility comparisons with their modifications, naturality in pure two-cells, and coherence with sequential composition and the specified pure pseudofunctor J. In particular this includes coherent centrality for pure state maps. These are the monoidal-action coherence data before introducing state machines, not axioms about a proposed state tricategory. The precise transferred action package is the one specified and constructed in the companion \emph{Strong pure actions after relative module completion}: original strong pure-strand action pseudofunctors, their action associators/unitors and bimodule modifications, and their compatibility with composition and $J$, lifted with extent-transported equivalence modules. In the intended completion these are the induced pure module tensor/action data. Merely having object-level tensor formulas is insufficient.

For fixed interface objects $A,B$, let $I(A,B)$ be the native dialgebra bicategory with objects $(X,f:JX\otimes A\to JX\otimes B)$ and arrows $(h,\alpha:h_B f\Rightarrow g h_A)$. A two-cell $\sigma:(h,\alpha)\Rightarrow(h',\alpha')$ is a state two-cell $\sigma:h\Rightarrow h'$ satisfying
\[
 (g\sigma_A)\alpha=\alpha'(\sigma_B f),
\]
with the actual associativity and action comparisons included. This becomes a three-cell in the proposed global structure. The forgetful hom pseudofunctor $U:I(A,B)\to S$ is faithful on these two-cells: they carry no extra datum beyond $\sigma$, although not every $\sigma$ satisfies the displayed equation.

\paragraph{Sequential assembly.} For $(X,f):A\to B$ and $(Y,g):B\to C$, use state $X\otimes Y$, run $f$ while holding $Y$, then $g$ while holding $X$, restoring the stated state order through pure action coherence. Denote this body $\operatorname{Seq}_{X,Y}(f,g)$. For native arrows $(h,\alpha)$ and $(k,\beta)$ the state map is $h\otimes k$. Paste the held copy of $\alpha$, then the held copy of $\beta$, using pure centrality to move the other state map past each body. This is the same typed pure-boundary rectangle as in the earlier all-intertwiner construction, now including the bicategorical action comparisons. On two-cells $(\sigma,\tau)$ use $\sigma\otimes\tau$.

To verify that this is a pseudofunctor $I(A,B)\times I(B,C)\to I(A,C)$, its unit and composition comparisons have underlying state cells given by the tensor pseudofunctor on $S$. Their native equations follow by expanding the two held-body pastes: the original square equations for $\alpha,\beta$, naturality of pure centrality at $\sigma,\tau$, and bicategorical interchange identify the two boundaries. Compatibility with vertical and horizontal composition of two-cells follows from the tensor pseudofunctor laws and the same naturality. Finally the pseudofunctor associativity and unit equations are detected by faithful $U$, once the comparison cells have been shown to satisfy their native equations. No inverse of an effectful sequential/parallel interchanger is used in this sequential construction.

\paragraph{Associator and units.} The two parenthesizations of three machines retain $(X\otimes Y)\otimes Z$ and $X\otimes(Y\otimes Z)$, and execute the same chronological body sequence $f,g,l$. Use the state associator equivalence and the corresponding action-body comparison as the associator native arrow. Its inverse and adjunction cells lift the specified state adjoint equivalence; the native equations are the action associativity equations. Its pseudonaturality on native arrows is the action associator's naturality, applied to the three held-body squares. The naturality modifications are the lifted state naturality two-cells, not assertions of strict equality. Unit machines have unit state and the identity body; unitors and their adjunction cells are constructed from the state/action unit data in the same way.

\paragraph{Held-body admissibility lemma.} Let $u,v$ be parallel structural state paths between two bracketings of the same retained state list, and let $\sigma:u\Rightarrow v$ be one of the specified structural coherence two-cells of the monoidal state/action doctrine, or a generated pasting of those cells. It is not an arbitrary state two-cell between these paths. Let $F,G$ be the corresponding held-body executions of the same chronological list of processes. The canonical action comparisons $\alpha_u:u_B F\Rightarrow G u_A$ and $\alpha_v:v_B F\Rightarrow G v_A$ satisfy
\[
 (G\sigma_A)\alpha_u=\alpha_v(\sigma_B F). \tag{*}
\]
For an empty process list this is action coherence on identities. For one held process it is the two-cell naturality of the relevant action associativity/unit/compatibility constraint. To append the next held process, paste that single-process equation with the induction hypothesis; bicategorical interchange identifies the two composite boundaries. Thus (*) holds for every chronological list and those specified structural cells. For naturality in arbitrary machine comparisons, their native comparison squares must be pasted into the argument; structural wire coherence does not replace those hypotheses. For example an arbitrary endomorphism $\sigma:1_X\Rightarrow1_X$ need not commute with a stateful transition and is not covered. Only structural wires are moved, not the effectful processes relative to each other. This proves admissibility, rather than inferring existence of a lift from faithfulness.

Apply (*) to the state associator's naturality cell on three held bodies, and to each state unitor's naturality cell on one held body. Apply it to the state pentagonator on four held bodies, and to the three monoidal-bicategory unit modifications $\lambda,\mu,\rho$ on the corresponding bodies with unit slots. Finally apply it to the units/counits of the associator/unitor adjoint equivalences on those same held bodies. In every case the native modification condition is exactly (*). Their modification naturality follows by the same equation with native input squares inserted and the action constraints' naturality. This specifies all the lifting checks; the later argument by faithful projection concerns only equations between these already constructed cells.

\paragraph{Tricategory convention.} We use the GPS/Gurski convention: hom bicategories, composition/identity pseudofunctors, pseudonatural adjoint equivalences for associativity and units, pentagonator $\pi$, and the three unit modifications $\lambda,\mu,\rho$, with their coherence equations. The data list is restated in Caviglia, \emph{2-Stacks over bisites}, Definition 1.1, \url{https://link.springer.com/article/10.1007/s40062-026-00405-z}, referring to Gurski Definition 3.1.2. The faithful-detection argument is for equations between already constructed modifications in this convention, not a replacement for their existence.

\paragraph{Higher coherence verification strategy.} The pentagonator and all three unit modifications must be actual native three-cells. Take their underlying state two-cells from the monoidal bicategory $S$. Their native square equations are the corresponding action pentagon and unit equations, evaluated on the identical chronological held bodies. This establishes their existence. All remaining equations between these three-cells are then detected by $U$, where they are precisely the monoidal bicategory coherence equations. This separates the essential lifting check from the subsequent faithfulness argument: faithfulness alone does not produce a missing native modification.

Under the specified action package, the preceding lifting audit yields a tricategory with interfaces as objects and $I(A,B)$ as hom bicategories. It does not claim a simultaneous strictification of the input or silently erase the state two-cells. If the state bicategory is locally discrete, the extra three-cells are identities and the construction recovers the previously established ordinary-state native bicategory.

\paragraph{Synchronization and next obligations.} In the derived module completion, the fixed-interface theorem supplies selected lax limits in these hom bicategories. To prove compatibility with horizontal assembly, use the earlier creation/preservation criterion only after the sequential assembly pseudofunctor has been verified: its underlying state tensor preserves the product-index cone presentations, and its held-body construction sends pseudo native arrows to pseudo arrows. This can address independent product-index assembly; same-index shared synchronization remains a different diagonal comparison.

Flattening requires another dimension check. Iterating the old Cat-valued state constructor is not immediately defined on this tricategory. A higher state constructor, its action on the relevant hom bicategories and structural transformations, and its unit/multiplication modifications must be specified before claiming a higher pseudomonad. The present horizontal construction therefore does not by itself supply global feedback evaluation, instantaneous trace, or general mixed shared-interface feedback.
\end{document}
